% Options for packages loaded elsewhere
% Options for packages loaded elsewhere
\PassOptionsToPackage{unicode}{hyperref}
\PassOptionsToPackage{hyphens}{url}
%
\documentclass[
  ignorenonframetext,
  aspectratio=169,
  spanish,
]{beamer}
\newif\ifbibliography
\usepackage{pgfpages}
\setbeamertemplate{caption}[numbered]
\setbeamertemplate{caption label separator}{: }
\setbeamercolor{caption name}{fg=normal text.fg}
\beamertemplatenavigationsymbolshorizontal
% remove section numbering
\setbeamertemplate{part page}{
  \centering
  \begin{beamercolorbox}[sep=16pt,center]{part title}
    \usebeamerfont{part title}\insertpart\par
  \end{beamercolorbox}
}
\setbeamertemplate{section page}{
  \centering
  \begin{beamercolorbox}[sep=12pt,center]{section title}
    \usebeamerfont{section title}\insertsection\par
  \end{beamercolorbox}
}
\setbeamertemplate{subsection page}{
  \centering
  \begin{beamercolorbox}[sep=8pt,center]{subsection title}
    \usebeamerfont{subsection title}\insertsubsection\par
  \end{beamercolorbox}
}
% Prevent slide breaks in the middle of a paragraph
\widowpenalties 1 10000
\raggedbottom
\AtBeginPart{
  \frame{\partpage}
}
\AtBeginSection{
  \ifbibliography
  \else
    \frame{\sectionpage}
  \fi
}
\AtBeginSubsection{
  \frame{\subsectionpage}
}
\usepackage{iftex}
\ifPDFTeX
  \usepackage[T1]{fontenc}
  \usepackage[utf8]{inputenc}
  \usepackage{textcomp} % provide euro and other symbols
\else % if luatex or xetex
  \usepackage{unicode-math} % this also loads fontspec
  \defaultfontfeatures{Scale=MatchLowercase}
  \defaultfontfeatures[\rmfamily]{Ligatures=TeX,Scale=1}
\fi
\usepackage{lmodern}

\usetheme[]{default}
\usefonttheme[]{default}
\ifPDFTeX\else
  % xetex/luatex font selection
\fi
% Use upquote if available, for straight quotes in verbatim environments
\IfFileExists{upquote.sty}{\usepackage{upquote}}{}
\IfFileExists{microtype.sty}{% use microtype if available
  \usepackage[]{microtype}
  \UseMicrotypeSet[protrusion]{basicmath} % disable protrusion for tt fonts
}{}
\makeatletter
\@ifundefined{KOMAClassName}{% if non-KOMA class
  \IfFileExists{parskip.sty}{%
    \usepackage{parskip}
  }{% else
    \setlength{\parindent}{0pt}
    \setlength{\parskip}{6pt plus 2pt minus 1pt}}
}{% if KOMA class
  \KOMAoptions{parskip=half}}
\makeatother

\usepackage{color}
\usepackage{fancyvrb}
\newcommand{\VerbBar}{|}
\newcommand{\VERB}{\Verb[commandchars=\\\{\}]}
\DefineVerbatimEnvironment{Highlighting}{Verbatim}{commandchars=\\\{\}}
% Add ',fontsize=\small' for more characters per line
\usepackage{framed}
\definecolor{shadecolor}{RGB}{241,243,245}
\newenvironment{Shaded}{\begin{snugshade}}{\end{snugshade}}
\newcommand{\AlertTok}[1]{\textcolor[rgb]{0.68,0.00,0.00}{#1}}
\newcommand{\AnnotationTok}[1]{\textcolor[rgb]{0.37,0.37,0.37}{#1}}
\newcommand{\AttributeTok}[1]{\textcolor[rgb]{0.40,0.45,0.13}{#1}}
\newcommand{\BaseNTok}[1]{\textcolor[rgb]{0.68,0.00,0.00}{#1}}
\newcommand{\BuiltInTok}[1]{\textcolor[rgb]{0.00,0.23,0.31}{#1}}
\newcommand{\CharTok}[1]{\textcolor[rgb]{0.13,0.47,0.30}{#1}}
\newcommand{\CommentTok}[1]{\textcolor[rgb]{0.37,0.37,0.37}{#1}}
\newcommand{\CommentVarTok}[1]{\textcolor[rgb]{0.37,0.37,0.37}{\textit{#1}}}
\newcommand{\ConstantTok}[1]{\textcolor[rgb]{0.56,0.35,0.01}{#1}}
\newcommand{\ControlFlowTok}[1]{\textcolor[rgb]{0.00,0.23,0.31}{\textbf{#1}}}
\newcommand{\DataTypeTok}[1]{\textcolor[rgb]{0.68,0.00,0.00}{#1}}
\newcommand{\DecValTok}[1]{\textcolor[rgb]{0.68,0.00,0.00}{#1}}
\newcommand{\DocumentationTok}[1]{\textcolor[rgb]{0.37,0.37,0.37}{\textit{#1}}}
\newcommand{\ErrorTok}[1]{\textcolor[rgb]{0.68,0.00,0.00}{#1}}
\newcommand{\ExtensionTok}[1]{\textcolor[rgb]{0.00,0.23,0.31}{#1}}
\newcommand{\FloatTok}[1]{\textcolor[rgb]{0.68,0.00,0.00}{#1}}
\newcommand{\FunctionTok}[1]{\textcolor[rgb]{0.28,0.35,0.67}{#1}}
\newcommand{\ImportTok}[1]{\textcolor[rgb]{0.00,0.46,0.62}{#1}}
\newcommand{\InformationTok}[1]{\textcolor[rgb]{0.37,0.37,0.37}{#1}}
\newcommand{\KeywordTok}[1]{\textcolor[rgb]{0.00,0.23,0.31}{\textbf{#1}}}
\newcommand{\NormalTok}[1]{\textcolor[rgb]{0.00,0.23,0.31}{#1}}
\newcommand{\OperatorTok}[1]{\textcolor[rgb]{0.37,0.37,0.37}{#1}}
\newcommand{\OtherTok}[1]{\textcolor[rgb]{0.00,0.23,0.31}{#1}}
\newcommand{\PreprocessorTok}[1]{\textcolor[rgb]{0.68,0.00,0.00}{#1}}
\newcommand{\RegionMarkerTok}[1]{\textcolor[rgb]{0.00,0.23,0.31}{#1}}
\newcommand{\SpecialCharTok}[1]{\textcolor[rgb]{0.37,0.37,0.37}{#1}}
\newcommand{\SpecialStringTok}[1]{\textcolor[rgb]{0.13,0.47,0.30}{#1}}
\newcommand{\StringTok}[1]{\textcolor[rgb]{0.13,0.47,0.30}{#1}}
\newcommand{\VariableTok}[1]{\textcolor[rgb]{0.07,0.07,0.07}{#1}}
\newcommand{\VerbatimStringTok}[1]{\textcolor[rgb]{0.13,0.47,0.30}{#1}}
\newcommand{\WarningTok}[1]{\textcolor[rgb]{0.37,0.37,0.37}{\textit{#1}}}

\usepackage{longtable,booktabs,array}
\usepackage{calc} % for calculating minipage widths
\usepackage{caption}
% Make caption package work with longtable
\makeatletter
\def\fnum@table{\tablename~\thetable}
\makeatother
\usepackage{graphicx}
\makeatletter
\newsavebox\pandoc@box
\newcommand*\pandocbounded[1]{% scales image to fit in text height/width
  \sbox\pandoc@box{#1}%
  \Gscale@div\@tempa{\textheight}{\dimexpr\ht\pandoc@box+\dp\pandoc@box\relax}%
  \Gscale@div\@tempb{\linewidth}{\wd\pandoc@box}%
  \ifdim\@tempb\p@<\@tempa\p@\let\@tempa\@tempb\fi% select the smaller of both
  \ifdim\@tempa\p@<\p@\scalebox{\@tempa}{\usebox\pandoc@box}%
  \else\usebox{\pandoc@box}%
  \fi%
}
% Set default figure placement to htbp
\def\fps@figure{htbp}
\makeatother



\ifLuaTeX
\usepackage[bidi=basic]{babel}
\else
\usepackage[bidi=default]{babel}
\fi
% get rid of language-specific shorthands (see #6817):
\let\LanguageShortHands\languageshorthands
\def\languageshorthands#1{}


\setlength{\emergencystretch}{3em} % prevent overfull lines

\providecommand{\tightlist}{%
  \setlength{\itemsep}{0pt}\setlength{\parskip}{0pt}}



 


\setbeamertemplate{headline}{} % Elimina la cabecera del tema Madrid
\setbeamertemplate{section in head/foot}{} % Limpia la numeración de secciones
\renewcommand{\contentsname}{Índice}
\renewcommand{\refname}{Referencias}
\renewcommand{\figurename}{Figura}
\renewcommand{\tablename}{Tabla}
\logo{%
  \begin{tabular}{c}
    \includegraphics[height=1.2cm]{EscudoU.png}
  \end{tabular}%
}
% --- CONFIGURACIÓN DE COLORES INSTITUCIONALES ---
\definecolor{UnicorGreen}{HTML}{024426}
\definecolor{UnicorText}{HTML}{F9F9F9}
\definecolor{UnicorWhite}{HTML}{FFFFFF}
\definecolor{UnicorGold}{HTML}{D1B061}

% Aplicación al fondo y texto general
\setbeamercolor{background canvas}{bg=UnicorGreen}
\setbeamercolor{normal text}{fg=UnicorText}

% Aplicación a títulos y cabeceras
%\setbeamercolor{title}{fg=UnicorWhite}
\setbeamercolor{frametitle}{fg=UnicorWhite}
\setbeamercolor{framesubtitle}{fg=UnicorWhite}

% Aplicación a viñetas y estructura
\setbeamercolor{structure}{fg=UnicorWhite}
\setbeamercolor{itemize item}{fg=UnicorWhite}
\setbeamercolor{itemize subitem}{fg=UnicorWhite}
\setbeamercolor{enumerate item}{fg=UnicorWhite}

% Forzar a LaTeX a repintar el texto general al iniciar el documento
\AtBeginDocument{\usebeamercolor[fg]{normal text}}

% Aplicación a hipervínculos
\hypersetup{colorlinks=true, urlcolor=UnicorGold, linkcolor=UnicorWhite}

% --- ACTIVAR ESTILO DE BLOQUES Y SOMBRAS (ESTILO MADRID) ---
\useinnertheme[shadow=true]{rounded}

% --- COLORES DEL RECUADRO DEL TÍTULO PRINCIPAL ---
% Fondo dorado con texto verde
\setbeamercolor{title}{bg=UnicorGold, fg=UnicorGreen}
\setbeamercolor{frametitle}{fg=UnicorWhite} % Mantiene el título de las diapositivas normal
\setbeamercolor{framesubtitle}{fg=UnicorWhite}

% --- COLORES PARA TEOREMAS, DEFINICIONES Y BLOQUES ---
% Título del bloque (Fondo dorado, texto verde)
\setbeamercolor{block title}{bg=UnicorGold, fg=UnicorGreen}
% Cuerpo del bloque (Fondo blanco, texto verde)
\setbeamercolor{block body}{bg=UnicorWhite, fg=UnicorGreen}

\AtBeginDocument{
  \author{%
    \begin{tabular}{c}
      Mario A. Morales Rivera \\
      \footnotesize \texttt{mamorales@correo.unicordoba.edu.co} \\
      \footnotesize Universidad de Córdoba
    \end{tabular}
  }
  \institute{}
}
% --- LIMPIAR DIAPOSITIVAS DE SECCIÓN (ELIMINAR "SECTION X") ---
\setbeamertemplate{section page}{
  \begin{center}
    \begin{beamercolorbox}[sep=12pt,center,rounded=true,shadow=true]{title}
      \usebeamerfont{title}\insertsection\par
    \end{beamercolorbox}
  \end{center}
}
\makeatletter
\@ifpackageloaded{tcolorbox}{}{\usepackage[skins,breakable]{tcolorbox}}
\@ifpackageloaded{fontawesome5}{}{\usepackage{fontawesome5}}
\definecolor{quarto-callout-color}{HTML}{909090}
\definecolor{quarto-callout-note-color}{HTML}{0758E5}
\definecolor{quarto-callout-important-color}{HTML}{CC1914}
\definecolor{quarto-callout-warning-color}{HTML}{EB9113}
\definecolor{quarto-callout-tip-color}{HTML}{00A047}
\definecolor{quarto-callout-caution-color}{HTML}{FC5300}
\definecolor{quarto-callout-color-frame}{HTML}{acacac}
\definecolor{quarto-callout-note-color-frame}{HTML}{4582ec}
\definecolor{quarto-callout-important-color-frame}{HTML}{d9534f}
\definecolor{quarto-callout-warning-color-frame}{HTML}{f0ad4e}
\definecolor{quarto-callout-tip-color-frame}{HTML}{02b875}
\definecolor{quarto-callout-caution-color-frame}{HTML}{fd7e14}
\makeatother
\makeatletter
\@ifpackageloaded{caption}{}{\usepackage{caption}}
\AtBeginDocument{%
\ifdefined\contentsname
  \renewcommand*\contentsname{Tabla de contenidos}
\else
  \newcommand\contentsname{Tabla de contenidos}
\fi
\ifdefined\listfigurename
  \renewcommand*\listfigurename{Listado de Figuras}
\else
  \newcommand\listfigurename{Listado de Figuras}
\fi
\ifdefined\listtablename
  \renewcommand*\listtablename{Listado de Tablas}
\else
  \newcommand\listtablename{Listado de Tablas}
\fi
\ifdefined\figurename
  \renewcommand*\figurename{Figura}
\else
  \newcommand\figurename{Figura}
\fi
\ifdefined\tablename
  \renewcommand*\tablename{Tabla}
\else
  \newcommand\tablename{Tabla}
\fi
}
\@ifpackageloaded{float}{}{\usepackage{float}}
\floatstyle{ruled}
\@ifundefined{c@chapter}{\newfloat{codelisting}{h}{lop}}{\newfloat{codelisting}{h}{lop}[chapter]}
\floatname{codelisting}{Listado}
\newcommand*\listoflistings{\listof{codelisting}{Listado de Listados}}
\makeatother
\makeatletter
\makeatother
\makeatletter
\@ifpackageloaded{caption}{}{\usepackage{caption}}
\@ifpackageloaded{subcaption}{}{\usepackage{subcaption}}
\makeatother
\makeatletter
\@ifpackageloaded{tikz}{}{\usepackage{tikz}}
\makeatother
        \newcommand*\circled[1]{\tikz[baseline=(char.base)]{
          \node[shape=circle,draw,inner sep=1pt] (char) {{\scriptsize#1}};}}  
                  

\usepackage{bookmark}
\IfFileExists{xurl.sty}{\usepackage{xurl}}{} % add URL line breaks if available
\urlstyle{same}
\hypersetup{
  pdftitle={Estadística con Python},
  pdfauthor={Mario Morales},
  pdflang={es},
  hidelinks,
  pdfcreator={LaTeX via pandoc}}


\title{Estadística con Python}
\subtitle{Tipos de Datos en Python - Calculadora}
\author{Mario Morales}
\date{}

\begin{document}
\frame{\titlepage}

\renewcommand*\contentsname{Tabla de contenidos}
\begin{frame}[allowframebreaks]
  \frametitle{Tabla de contenidos}
  \setcounter{tocdepth}{3}
  \tableofcontents
\end{frame}
\setcounter{tocdepth}{3}
\tableofcontents
}

\section{Objetivos de la Sesión}\label{objetivos-de-la-sesiuxf3n}

\begin{frame}{Objetivos de la Sesión}
\begin{itemize}
\tightlist
\item
  Comprender la naturaleza y utilidad de los diferentes tipos de datos
  en Python.
\item
  Diferenciar entre datos numéricos (enteros y flotantes), cadenas de
  texto y valores lógicos.
\item
  Reconocer estructuras de datos introductorias relevantes para el
  análisis de datos.
\item
  Aplicar funciones nativas de Python para verificar y convertir tipos
  de datos.
\end{itemize}
\end{frame}

\section{Tipos de datos: importancia}\label{tipos-de-datos-importancia}

\begin{frame}{Tipos de datos: importancia}
\begin{itemize}
\item
  En estadística y programación, la forma en que el computador almacena
  la información dicta qué operaciones matemáticas o lógicas podemos
  realizar.
\item
  No podemos calcular la media de un texto, ni concatenar
  matemáticamente dos números lógicos sin una conversión previa.
\item
  Python es un lenguaje de tipado dinámico, lo que significa que el tipo
  de dato se asigna automáticamente durante la ejecución.
\end{itemize}
\end{frame}

\section{Datos Numéricos}\label{datos-numuxe9ricos}

\begin{frame}{Datos Numéricos}
Python maneja los números principalmente de dos maneras fundamentales
para nuestros cálculos.
\end{frame}

\begin{frame}[fragile]{Enteros (\texttt{int})}
\phantomsection\label{enteros-int}
Representan números completos sin parte decimal. Son útiles para contar
frecuencias, iteraciones, o tamaños de muestra (\(n\)).

\begin{Shaded}
\begin{Highlighting}[]
\CommentTok{\# Ejemplo de Integer}
\NormalTok{tamaño\_muestra }\OperatorTok{=} \DecValTok{150}
\BuiltInTok{print}\NormalTok{(tamaño\_muestra)}
\end{Highlighting}
\end{Shaded}

\begin{verbatim}
150
\end{verbatim}

\begin{Shaded}
\begin{Highlighting}[]
\BuiltInTok{print}\NormalTok{(}\BuiltInTok{type}\NormalTok{(tamaño\_muestra))}
\end{Highlighting}
\end{Shaded}

\begin{verbatim}
<class 'int'>
\end{verbatim}
\end{frame}

\begin{frame}[fragile]{Flotantes (\texttt{floats})}
\phantomsection\label{flotantes-floats}
Representan números reales (con parte decimal). Son el estándar para
variables continuas, probabilidades, promedios y estadísticos de prueba.

\begin{Shaded}
\begin{Highlighting}[]
\CommentTok{\# Ejemplo de Float}
\NormalTok{media\_muestral }\OperatorTok{=} \FloatTok{45.89}
\NormalTok{p\_valor }\OperatorTok{=} \FloatTok{0.035}
\BuiltInTok{print}\NormalTok{(media\_muestral)}
\end{Highlighting}
\end{Shaded}

\begin{verbatim}
45.89
\end{verbatim}

\begin{Shaded}
\begin{Highlighting}[]
\BuiltInTok{print}\NormalTok{(}\BuiltInTok{type}\NormalTok{(media\_muestral))}
\end{Highlighting}
\end{Shaded}

\begin{verbatim}
<class 'float'>
\end{verbatim}
\end{frame}

\begin{frame}[fragile]{Datos de Texto o Caracteres (\texttt{str})}
\phantomsection\label{datos-de-texto-o-caracteres-str}
Conocidos como \texttt{strings}, almacenan texto. En el modelado
estadístico, a menudo usamos \texttt{strings} para etiquetas, nombres de
variables categóricas o descripciones cualitativas. Se delimitan con
comillas simples o dobles.

\begin{Shaded}
\begin{Highlighting}[]
\CommentTok{\# Ejemplo de String}
\NormalTok{categoria\_tratamiento }\OperatorTok{=} \StringTok{"Tratamiento A"}
\BuiltInTok{print}\NormalTok{(categoria\_tratamiento)}
\end{Highlighting}
\end{Shaded}

\begin{verbatim}
Tratamiento A
\end{verbatim}

\begin{Shaded}
\begin{Highlighting}[]
\BuiltInTok{print}\NormalTok{(}\BuiltInTok{type}\NormalTok{(categoria\_tratamiento))}
\end{Highlighting}
\end{Shaded}

\begin{verbatim}
<class 'str'>
\end{verbatim}
\end{frame}

\begin{frame}[fragile]{Datos Lógicos o Booleanos (\texttt{bool})}
\phantomsection\label{datos-luxf3gicos-o-booleanos-bool}
\begin{itemize}
\item
  Solo pueden tomar dos valores: \texttt{True} (Verdadero) o
  \texttt{False} (Falso).
\item
  Son la base del control de flujo y la evaluación de condiciones (por
  ejemplo, si un valor supera un umbral crítico).
\item
  En el fondo, Python evalúa True como \(1\) y False como \(0\).
\end{itemize}

\begin{Shaded}
\begin{Highlighting}[]
\CommentTok{\# Ejemplo de Booleano}
\NormalTok{hipotesis\_rechazada }\OperatorTok{=} \VariableTok{True}
\BuiltInTok{print}\NormalTok{(hipotesis\_rechazada)}
\end{Highlighting}
\end{Shaded}

\begin{verbatim}
True
\end{verbatim}

\begin{Shaded}
\begin{Highlighting}[]
\BuiltInTok{print}\NormalTok{(}\BuiltInTok{type}\NormalTok{(hipotesis\_rechazada))}
\end{Highlighting}
\end{Shaded}

\begin{verbatim}
<class 'bool'>
\end{verbatim}

\begin{Shaded}
\begin{Highlighting}[]
\CommentTok{\# Evaluación de una condición que retorna un booleano}
\NormalTok{es\_significativo }\OperatorTok{=}\NormalTok{ p\_valor }\OperatorTok{\textless{}} \FloatTok{0.05}
\BuiltInTok{print}\NormalTok{(}\StringTok{"¿Es estadísticamente significativo?:"}\NormalTok{, es\_significativo)}
\end{Highlighting}
\end{Shaded}

\begin{verbatim}
¿Es estadísticamente significativo?: True
\end{verbatim}
\end{frame}

\begin{frame}[fragile]{Ejemplo de \texttt{bool}}
\phantomsection\label{ejemplo-de-bool}
\begin{Shaded}
\begin{Highlighting}[]
\CommentTok{\# Evaluación de una condición que retorna un booleano}
\NormalTok{es\_significativo }\OperatorTok{=}\NormalTok{ p\_valor }\OperatorTok{\textless{}} \FloatTok{0.05}
\BuiltInTok{print}\NormalTok{(}\StringTok{"¿Es estadísticamente significativo?:"}\NormalTok{, es\_significativo)}
\end{Highlighting}
\end{Shaded}

\begin{verbatim}
¿Es estadísticamente significativo?: True
\end{verbatim}
\end{frame}

\begin{frame}[fragile]{Datos Especiales: Ausencia de Valor (NoneType)}
\phantomsection\label{datos-especiales-ausencia-de-valor-nonetype}
\begin{itemize}
\item
  En la limpieza de datos, es común encontrar valores perdidos
  (\texttt{Missing\ values}).
\item
  En Python puro, esto suele representarse con \texttt{None}.
\end{itemize}

\begin{Shaded}
\begin{Highlighting}[]
\CommentTok{\# Ejemplo de None}
\NormalTok{dato\_faltante }\OperatorTok{=} \VariableTok{None}
\BuiltInTok{print}\NormalTok{(}\BuiltInTok{type}\NormalTok{(dato\_faltante))}
\end{Highlighting}
\end{Shaded}

\begin{verbatim}
<class 'NoneType'>
\end{verbatim}
\end{frame}

\begin{frame}[fragile]{Conversión de Tipos (\texttt{Casting})}
\phantomsection\label{conversiuxf3n-de-tipos-casting}
A veces necesitamos transformar un tipo de dato en otro, por ejemplo,
convertir un texto numérico recolectado de un formulario web a un
entero.

\begin{Shaded}
\begin{Highlighting}[]
\CommentTok{\# De string a entero}
\NormalTok{texto\_numero }\OperatorTok{=} \StringTok{"120"}
\NormalTok{numero\_real }\OperatorTok{=} \BuiltInTok{int}\NormalTok{(texto\_numero)}
\BuiltInTok{print}\NormalTok{(}\BuiltInTok{type}\NormalTok{(numero\_real))}
\end{Highlighting}
\end{Shaded}

\begin{verbatim}
<class 'int'>
\end{verbatim}

\begin{Shaded}
\begin{Highlighting}[]
\CommentTok{\# De float a int (¡Cuidado: se pierde la parte decimal!)}
\NormalTok{valor\_estimado }\OperatorTok{=} \FloatTok{4.95}
\NormalTok{entero\_truncado }\OperatorTok{=} \BuiltInTok{int}\NormalTok{(valor\_estimado)}
\BuiltInTok{print}\NormalTok{(entero\_truncado)}
\end{Highlighting}
\end{Shaded}

\begin{verbatim}
4
\end{verbatim}
\end{frame}

\begin{frame}[fragile]{Ejercicio Práctico en Clase}
\phantomsection\label{ejercicio-pruxe1ctico-en-clase}
\begin{itemize}
\item
  Crear una variable llamada \texttt{edad} y asígnenle su edad en años
  (¿Qué tipo de dato es?).
\item
  Crear una variable estatura en metros (ej. 1.75).
\item
  Crear una variable \texttt{nombre\_estudiante} con su nombre completo.
\item
  Evaluar de forma lógica (Booleana) si la edad es mayor a 18 y guardar
  el resultado en una variable llamada \texttt{es\_mayor\_edad}.
\item
  Usen la función \texttt{print()} para mostrar todas las variables y la
  función \texttt{type()} para verificar que usaron el tipo correcto.
\end{itemize}
\end{frame}

\begin{frame}[fragile]{Ejemplo de aplicación}
\phantomsection\label{ejemplo-de-aplicaciuxf3n}
\begin{itemize}
\item
  Escribir un programa es \texttt{python} que capture dos números
  \texttt{a} y \texttt{b} y regrese la suma \texttt{a+b}
\item
  \emph{Ayuda:} para capturar un valor desde el teclado use la función
  \texttt{input()}
\end{itemize}
\end{frame}

\section{Python como calculadora}\label{python-como-calculadora}

\begin{frame}[fragile]{Operaciones I}
\phantomsection\label{operaciones-i}
\begin{Shaded}
\begin{Highlighting}[]
\CommentTok{\# Suma }
\DecValTok{2}\OperatorTok{+}\DecValTok{2}
\end{Highlighting}
\end{Shaded}

\begin{verbatim}
4
\end{verbatim}

\begin{Shaded}
\begin{Highlighting}[]
\CommentTok{\# Resta }
\DecValTok{5}\OperatorTok{{-}}\DecValTok{3} 
\end{Highlighting}
\end{Shaded}

\begin{verbatim}
2
\end{verbatim}

\begin{Shaded}
\begin{Highlighting}[]
\CommentTok{\# Producto }
\DecValTok{3}\OperatorTok{*}\DecValTok{5} 
\end{Highlighting}
\end{Shaded}

\begin{verbatim}
15
\end{verbatim}
\end{frame}

\begin{frame}[fragile]{Operaciones II}
\phantomsection\label{operaciones-ii}
\begin{Shaded}
\begin{Highlighting}[]
\CommentTok{\#\# Cuidado }
\DecValTok{2}\OperatorTok{\^{}}\DecValTok{2} \CommentTok{\# no funciona como potencia     }
\end{Highlighting}
\end{Shaded}

\begin{verbatim}
0
\end{verbatim}

\begin{Shaded}
\begin{Highlighting}[]
\CommentTok{\#\#\# para la potenciación en python potencia se debe usar **}
\BuiltInTok{print}\NormalTok{(}\DecValTok{2}\OperatorTok{**}\DecValTok{2}\NormalTok{)}
\end{Highlighting}
\end{Shaded}

\begin{verbatim}
4
\end{verbatim}

\begin{Shaded}
\begin{Highlighting}[]
\DecValTok{9}\OperatorTok{**}\NormalTok{(}\DecValTok{1}\OperatorTok{/}\DecValTok{2}\NormalTok{) }
\end{Highlighting}
\end{Shaded}

\begin{verbatim}
3.0
\end{verbatim}
\end{frame}

\begin{frame}[fragile]{Operaciones III}
\phantomsection\label{operaciones-iii}
\begin{Shaded}
\begin{Highlighting}[]
\CommentTok{\# El orden de las operaciones se altera con paréntesis }

\NormalTok{(}\DecValTok{4}\OperatorTok{{-}}\DecValTok{2}\NormalTok{)}\OperatorTok{*}\DecValTok{3}\OperatorTok{**}\DecValTok{2} 
\end{Highlighting}
\end{Shaded}

\begin{verbatim}
18
\end{verbatim}

\begin{Shaded}
\begin{Highlighting}[]
\CommentTok{\# Residuo de la división}
\DecValTok{7}\OperatorTok{\%}\DecValTok{4}
\end{Highlighting}
\end{Shaded}

\begin{verbatim}
3
\end{verbatim}

\begin{Shaded}
\begin{Highlighting}[]
\CommentTok{\# Parte entera de la división }
\DecValTok{7}\OperatorTok{//}\DecValTok{4}
\end{Highlighting}
\end{Shaded}

\begin{verbatim}
1
\end{verbatim}
\end{frame}

\begin{frame}[fragile]{Funciones I}
\phantomsection\label{funciones-i}
Para tener disponibles funciones matemáticas es necesario importar el
paquete \texttt{math}

\begin{Shaded}
\begin{Highlighting}[]
\CommentTok{\# Raíz cuadrada }
\ImportTok{from}\NormalTok{ math }\ImportTok{import}\NormalTok{ sqrt}
\NormalTok{sqrt(}\DecValTok{9}\NormalTok{)}
\end{Highlighting}
\end{Shaded}

\begin{verbatim}
3.0
\end{verbatim}

\begin{Shaded}
\begin{Highlighting}[]
\CommentTok{\# exponencial }
\ImportTok{from}\NormalTok{ math }\ImportTok{import}\NormalTok{ exp}
\NormalTok{exp(}\DecValTok{1}\NormalTok{)}
\end{Highlighting}
\end{Shaded}

\begin{verbatim}
2.718281828459045
\end{verbatim}

\begin{Shaded}
\begin{Highlighting}[]
\CommentTok{\# logaritmo en base 10 }
\ImportTok{from}\NormalTok{ math }\ImportTok{import}\NormalTok{ log10}
\NormalTok{log10(}\DecValTok{100}\NormalTok{) }
\end{Highlighting}
\end{Shaded}

\begin{verbatim}
2.0
\end{verbatim}
\end{frame}

\begin{frame}[fragile]{Funciones II}
\phantomsection\label{funciones-ii}
\begin{Shaded}
\begin{Highlighting}[]
\ImportTok{from}\NormalTok{ math }\ImportTok{import}\NormalTok{ log}
\CommentTok{\# logaritmo en base 10 }
\NormalTok{log(}\DecValTok{100}\NormalTok{,}\DecValTok{10}\NormalTok{)}
\end{Highlighting}
\end{Shaded}

\begin{verbatim}
2.0
\end{verbatim}

\begin{Shaded}
\begin{Highlighting}[]
\CommentTok{\# logaritmo en base 2 }
\NormalTok{log(}\DecValTok{8}\NormalTok{,}\DecValTok{2}\NormalTok{)}
\end{Highlighting}
\end{Shaded}

\begin{verbatim}
3.0
\end{verbatim}
\end{frame}

\begin{frame}{Funciones trigonométricas}
\phantomsection\label{funciones-trigonomuxe9tricas}
Usando Python, calcular

\begin{itemize}
\tightlist
\item
  \(\sin(\pi/2)\)
\item
  \(\sin(\pi)\)
\item
  \(\sin(2\times \pi)\)
\item
  \(\cos(\pi/2)\)
\item
  \(\cos(\pi)\)
\item
  \(\cos(2\times \pi)\)
\item
  \(\tan(\pi/4)\)
\end{itemize}
\end{frame}

\begin{frame}[fragile]{Solución}
\phantomsection\label{soluciuxf3n}
\begin{Shaded}
\begin{Highlighting}[]
\ImportTok{from}\NormalTok{ math }\ImportTok{import}\NormalTok{ sin, cos , pi }
\end{Highlighting}
\end{Shaded}
\end{frame}

\begin{frame}[fragile]{Tuplas}
\phantomsection\label{tuplas}
Las tuplas sirven para agrupar objetos de diferente tipo, ellas estan
\emph{indexadas} y son \textbf{inmutables}, es decir, no se pueden
modificar (agregar, quitar, cambiar entradas) una vez creadas.

\begin{Shaded}
\begin{Highlighting}[]
\NormalTok{mytupla0}\OperatorTok{=}\NormalTok{(}\StringTok{\textquotesingle{}abc\textquotesingle{}}\NormalTok{,}\DecValTok{4}\NormalTok{,}\VariableTok{True}\NormalTok{,}\DecValTok{5}\NormalTok{,}\VariableTok{False}\NormalTok{,}\FloatTok{2.5}\NormalTok{)}
\NormalTok{mytupla}\OperatorTok{=}\NormalTok{(}\DecValTok{1}\NormalTok{,}\FloatTok{0.7}\NormalTok{,mytupla0,}\StringTok{"Hola"}\NormalTok{) }
\BuiltInTok{print}\NormalTok{(mytupla0)}
\end{Highlighting}
\end{Shaded}

\begin{verbatim}
('abc', 4, True, 5, False, 2.5)
\end{verbatim}

\begin{Shaded}
\begin{Highlighting}[]
\NormalTok{mytupla}
\end{Highlighting}
\end{Shaded}

\begin{verbatim}
(1, 0.7, ('abc', 4, True, 5, False, 2.5), 'Hola')
\end{verbatim}
\end{frame}

\begin{frame}{Datos soportados por las tuplas}
\phantomsection\label{datos-soportados-por-las-tuplas}
\begin{tcolorbox}[enhanced jigsaw, colbacktitle=quarto-callout-important-color!10!white, colframe=quarto-callout-important-color-frame, arc=.35mm, opacityback=0, toptitle=1mm, left=2mm, breakable, title=\textcolor{quarto-callout-important-color}{\faExclamation}\hspace{0.5em}{Importante}, bottomtitle=1mm, titlerule=0mm, rightrule=.15mm, coltitle=black, bottomrule=.15mm, toprule=.15mm, colback=white, leftrule=.75mm, opacitybacktitle=0.6]

Note que una tupla puede contener distintos tipos de datos (string,
enteros, lógicos y floats), incluso otras tuplas y otros tipos de
contenedores

\end{tcolorbox}
\end{frame}

\begin{frame}[fragile]{Acceso a datos en una tupla}
\phantomsection\label{acceso-a-datos-en-una-tupla}
Para acceder a los elementos de una tupla se usa
\texttt{tupla{[}indice{]}} donde \emph{indice} es un número entero de
cero hasta la longitud de la tupla menos uno \(n-1\)

\phantomsection\label{annotated-cell-31}%
\begin{Shaded}
\begin{Highlighting}[]
\NormalTok{mytupla[}\DecValTok{0}\NormalTok{] }\hspace*{\fill}\NormalTok{\circled{1}}
\end{Highlighting}
\end{Shaded}

\begin{description}
\tightlist
\item[\circled{1}]
Primer elemento de la tupla.
\end{description}

\begin{verbatim}
1
\end{verbatim}
\end{frame}

\begin{frame}[fragile]{Ejecute I}
\phantomsection\label{ejecute-i}
\begin{Shaded}
\begin{Highlighting}[]
\NormalTok{mytupla[}\DecValTok{4}\NormalTok{]}
\NormalTok{mytupla[}\DecValTok{2}\NormalTok{]}\OperatorTok{=}\DecValTok{0} 
\NormalTok{mytupla[}\DecValTok{2}\NormalTok{]}\OperatorTok{=}\StringTok{\textquotesingle{}xyz\textquotesingle{}}
\end{Highlighting}
\end{Shaded}
\end{frame}

\begin{frame}[fragile]{Ejecute II}
\phantomsection\label{ejecute-ii}
\begin{Shaded}
\begin{Highlighting}[]
\NormalTok{mytupla[}\DecValTok{4}\NormalTok{] }\CommentTok{\# IndexError: tuple index out of range}
\CommentTok{\# no corre, a eso se refiere con inmutable }
\NormalTok{mytupla[}\DecValTok{2}\NormalTok{]}\OperatorTok{=}\DecValTok{0} \CommentTok{\#TypeError: \textquotesingle{}tuple\textquotesingle{} object does not support item assignment}
\NormalTok{mytupla[}\DecValTok{2}\NormalTok{]}\OperatorTok{=}\StringTok{\textquotesingle{}xyz\textquotesingle{}}
\end{Highlighting}
\end{Shaded}
\end{frame}

\begin{frame}[fragile]{Observe}
\phantomsection\label{observe}
\begin{Shaded}
\begin{Highlighting}[]
\NormalTok{t1}\OperatorTok{=}\NormalTok{(}\DecValTok{3}\NormalTok{,}\DecValTok{4}\NormalTok{) }
\BuiltInTok{print}\NormalTok{(}\StringTok{\textquotesingle{}t1= \textquotesingle{}}\NormalTok{,t1,}\StringTok{\textquotesingle{}id de t1 antes \textquotesingle{}}\NormalTok{,}\BuiltInTok{id}\NormalTok{(t1))}
\end{Highlighting}
\end{Shaded}

\begin{verbatim}
t1=  (3, 4) id de t1 antes  4778509312
\end{verbatim}

\begin{Shaded}
\begin{Highlighting}[]
\NormalTok{t2}\OperatorTok{=}\NormalTok{(}\DecValTok{5}\NormalTok{,}\DecValTok{6}\NormalTok{)}
\BuiltInTok{print}\NormalTok{(}\StringTok{\textquotesingle{}t2= \textquotesingle{}}\NormalTok{,t2,}\StringTok{\textquotesingle{}id de t2 antes \textquotesingle{}}\NormalTok{,}\BuiltInTok{id}\NormalTok{(t2))}
\end{Highlighting}
\end{Shaded}

\begin{verbatim}
t2=  (5, 6) id de t2 antes  4778496832
\end{verbatim}

\begin{Shaded}
\begin{Highlighting}[]
\BuiltInTok{print}\NormalTok{(t2)}
\end{Highlighting}
\end{Shaded}

\begin{verbatim}
(5, 6)
\end{verbatim}

\begin{Shaded}
\begin{Highlighting}[]
\NormalTok{t1}\OperatorTok{=}\NormalTok{t1}\OperatorTok{+}\NormalTok{t2}
\BuiltInTok{print}\NormalTok{(}\StringTok{\textquotesingle{}t1= \textquotesingle{}}\NormalTok{, t1,}\StringTok{\textquotesingle{}id de t1 despues \textquotesingle{}}\NormalTok{, }\BuiltInTok{id}\NormalTok{(t1))}
\end{Highlighting}
\end{Shaded}

\begin{verbatim}
t1=  (3, 4, 5, 6) id de t1 despues  4778839392
\end{verbatim}

\begin{Shaded}
\begin{Highlighting}[]
\BuiltInTok{print}\NormalTok{(}\StringTok{\textquotesingle{}t2= \textquotesingle{}}\NormalTok{, t2, }\StringTok{\textquotesingle{}id de t2 despues \textquotesingle{}}\NormalTok{, }\BuiltInTok{id}\NormalTok{(t2))}
\end{Highlighting}
\end{Shaded}

\begin{verbatim}
t2=  (5, 6) id de t2 despues  4778496832
\end{verbatim}
\end{frame}

\begin{frame}{Ventajas de la inmutabilidad}
\phantomsection\label{ventajas-de-la-inmutabilidad}
\begin{itemize}
\item
  Seguridad: Como los datos no pueden ser modificados accidentalmente,
  es más seguro usar tuplas cuando se quiere proteger la integridad de
  los datos.
\item
  Uso como claves en diccionarios: Las tuplas, al ser inmutables, pueden
  ser utilizadas como claves en diccionarios, algo que no es posible con
  las listas.
\item
  Optimización: Las tuplas pueden ser más eficientes en términos de
  rendimiento y uso de memoria en comparación con las listas. La
  inmutabilidad hace que las tuplas sean útiles en situaciones donde se
  requiere un conjunto de datos constante y seguro.
\end{itemize}
\end{frame}

\begin{frame}[fragile]{Función \texttt{type}}
\phantomsection\label{funciuxf3n-type}
Para averiguar la \textbf{clase} de un objeto usamos la función
\texttt{type()}

\begin{Shaded}
\begin{Highlighting}[]
\BuiltInTok{type}\NormalTok{(mytupla) }
\end{Highlighting}
\end{Shaded}

\begin{verbatim}
<class 'tuple'>
\end{verbatim}
\end{frame}

\begin{frame}[fragile]{Averiguar si un objeto es una tupla}
\phantomsection\label{averiguar-si-un-objeto-es-una-tupla}
\begin{itemize}
\item
  En algunas aplicaciones es necesario verificar si un objeto pertenece
  o no a una clase.
\item
  Por ejemplo, para dirigir el flujo de ejecución del código de
  programación de una función: si el objeto entregado es una tupla,
  ejecute ciertas líneas de código, de lo contrario ejecute otras
  líneas.
\item
  La función \texttt{isinstance(objeto,tipo)} regresa \texttt{True} si
  la clase del \texttt{objeto} entregado corresponde a \texttt{tipo}, en
  caso contrario regresa \texttt{False}.
\item
  Otra forma es mediante la función \texttt{type}
\end{itemize}
\end{frame}

\begin{frame}[fragile]{Función \texttt{isinstance}}
\phantomsection\label{funciuxf3n-isinstance}
\begin{Shaded}
\begin{Highlighting}[]
\BuiltInTok{isinstance}\NormalTok{(mytupla, }\BuiltInTok{tuple}\NormalTok{)}
\end{Highlighting}
\end{Shaded}

\begin{verbatim}
True
\end{verbatim}

\begin{Shaded}
\begin{Highlighting}[]
\BuiltInTok{isinstance}\NormalTok{(}\StringTok{\textquotesingle{}abcd\textquotesingle{}}\NormalTok{, }\BuiltInTok{tuple}\NormalTok{)}
\end{Highlighting}
\end{Shaded}

\begin{verbatim}
False
\end{verbatim}

\begin{Shaded}
\begin{Highlighting}[]
\BuiltInTok{isinstance}\NormalTok{(}\StringTok{\textquotesingle{}abcd\textquotesingle{}}\NormalTok{, }\BuiltInTok{str}\NormalTok{)}
\end{Highlighting}
\end{Shaded}

\begin{verbatim}
True
\end{verbatim}
\end{frame}

\begin{frame}[fragile]{Función \texttt{isinstance}}
\phantomsection\label{funciuxf3n-isinstance-1}
\begin{Shaded}
\begin{Highlighting}[]
\BuiltInTok{type}\NormalTok{(mytupla)}
\end{Highlighting}
\end{Shaded}

\begin{verbatim}
<class 'tuple'>
\end{verbatim}

\begin{Shaded}
\begin{Highlighting}[]
\BuiltInTok{type}\NormalTok{(mytupla) }\KeywordTok{is} \BuiltInTok{tuple}
\end{Highlighting}
\end{Shaded}

\begin{verbatim}
True
\end{verbatim}

\begin{Shaded}
\begin{Highlighting}[]
\BuiltInTok{type}\NormalTok{(}\DecValTok{4}\NormalTok{) }\KeywordTok{is} \BuiltInTok{tuple}
\end{Highlighting}
\end{Shaded}

\begin{verbatim}
False
\end{verbatim}

\begin{Shaded}
\begin{Highlighting}[]
\BuiltInTok{type}\NormalTok{(}\DecValTok{4}\NormalTok{) }\KeywordTok{is} \BuiltInTok{int}
\end{Highlighting}
\end{Shaded}

\begin{verbatim}
True
\end{verbatim}
\end{frame}

\begin{frame}[fragile]{Otras formas de crear tuplas}
\phantomsection\label{otras-formas-de-crear-tuplas}
Otra forma de crear tuplas es mediante la siguiente orden, con la
segunda linea se imprime el tipo de objeto que es \texttt{t3}, con el
fin de verificar que efectivamente se ha creado una tupla.

\begin{Shaded}
\begin{Highlighting}[]
\NormalTok{t3}\OperatorTok{=}\DecValTok{10}\NormalTok{,}\DecValTok{20} 
\BuiltInTok{type}\NormalTok{(t3)  }
\end{Highlighting}
\end{Shaded}

\begin{verbatim}
<class 'tuple'>
\end{verbatim}

Esta forma se conoce como creación de tuplas mediante \emph{empaquetado}
\end{frame}

\begin{frame}[fragile]{Desempaquetado de tuplas}
\phantomsection\label{desempaquetado-de-tuplas}
También podemos recuperar los elementos de una tupla mediante el
\emph{desempaquetado}\index{Tuplas!desempaquetado}, de la siguiente
manera

\begin{Shaded}
\begin{Highlighting}[]
\NormalTok{x,y}\OperatorTok{=}\NormalTok{t3}
\BuiltInTok{print}\NormalTok{(}\StringTok{"x= "}\NormalTok{,x)}
\BuiltInTok{print}\NormalTok{(}\StringTok{"y= "}\NormalTok{,y)}
\end{Highlighting}
\end{Shaded}

el primer elemento de la tupla \texttt{t3} se asigna al objeto
\texttt{x} y el segundo a \texttt{y} luego se imprimen sus valores.
\end{frame}

\begin{frame}[fragile]{Tuplas especiales}
\phantomsection\label{tuplas-especiales}
Se pueden crear tuplas vacías y tuplas que contiene un solo valor, estas
últimas llamadas tuplas unitarias, como se muestra a continuación

\begin{Shaded}
\begin{Highlighting}[]
\NormalTok{tupvacia}\OperatorTok{=}\NormalTok{() }\CommentTok{\# crea una tupla vacía }
\NormalTok{tupunit}\OperatorTok{=}\NormalTok{(}\DecValTok{4}\NormalTok{,) }\CommentTok{\# crea una tupla con un único elemento }
\BuiltInTok{print}\NormalTok{(tupvacia,}\BuiltInTok{type}\NormalTok{(tupvacia),tupunit,}\BuiltInTok{type}\NormalTok{(tupunit))}
\end{Highlighting}
\end{Shaded}

\begin{verbatim}
() <class 'tuple'> (4,) <class 'tuple'>
\end{verbatim}

note el uso de la coma después del elemento de la tupla unitaria, eso
debe ser así porque de lo contrario no se estaría creando una tupla,
sino un objeto de tipo entero (\texttt{int})

\begin{Shaded}
\begin{Highlighting}[]
\NormalTok{x}\OperatorTok{=}\NormalTok{(}\DecValTok{4}\NormalTok{) }\CommentTok{\# cuidado!!! x es un entero, no una tupla  }
\BuiltInTok{print}\NormalTok{(x,}\BuiltInTok{type}\NormalTok{(x))}
\end{Highlighting}
\end{Shaded}

\begin{verbatim}
4 <class 'int'>
\end{verbatim}
\end{frame}

\begin{frame}[fragile]{Indexado de tuplas.}
\phantomsection\label{indexado-de-tuplas.}
\index{Tuplas!Indexado de}

Se puede crear una tupla a partir de otra mediante órdenes de la forma
\texttt{tupla{[}i:j{]}} esto se conoce como indexación, veamos algunos
ejemplos:

\begin{Shaded}
\begin{Highlighting}[]
\NormalTok{a}\OperatorTok{=}\NormalTok{(}\DecValTok{3}\NormalTok{,}\DecValTok{2}\NormalTok{,}\DecValTok{4}\NormalTok{,}\DecValTok{1}\NormalTok{,}\DecValTok{6}\NormalTok{,}\DecValTok{4}\NormalTok{) }\CommentTok{\# se crea la tupla a }
\NormalTok{b}\OperatorTok{=}\NormalTok{a[}\DecValTok{1}\NormalTok{:}\DecValTok{4}\NormalTok{] }\CommentTok{\# tupla b, creada a partir de a }
\BuiltInTok{print}\NormalTok{(b,}\BuiltInTok{type}\NormalTok{(b))}
\end{Highlighting}
\end{Shaded}

\begin{verbatim}
(2, 4, 1) <class 'tuple'>
\end{verbatim}

Note en la salida anterior que la tupla \texttt{b} inicia en \texttt{2},
y no en \texttt{3}, que es el primer elemento de la tupla \texttt{a}
¿por qué?

\begin{Shaded}
\begin{Highlighting}[]
\NormalTok{a[:}\DecValTok{4}\NormalTok{]    }\CommentTok{\# los cuatro primeros elementos de la tupla a}
\NormalTok{a[}\DecValTok{2}\NormalTok{:]    }\CommentTok{\# desde el tercer elemento hasta el final }
\NormalTok{a[::}\OperatorTok{{-}}\DecValTok{1}\NormalTok{]  }\CommentTok{\# la tupla en reversa  }
\end{Highlighting}
\end{Shaded}
\end{frame}

\section{Funciones y métodos aplicables a
tuplas}\label{sec-funcionesAplicablesATuplas}

\begin{frame}{Funciones y métodos aplicables a tuplas}
\index{Tuplas!Función!\texttt{len()}}\index{Tuplas!Método!\texttt{.index}}\index{Tuplas!Método!\texttt{.count()}}
\end{frame}

\begin{frame}[fragile]{La función \texttt{len()}}
\phantomsection\label{la-funciuxf3n-len}
\begin{itemize}
\tightlist
\item
  La función \texttt{len()} regresa el número de elementos que contiene
  una tupla
\end{itemize}

\begin{Shaded}
\begin{Highlighting}[]
\BuiltInTok{print}\NormalTok{(}\BuiltInTok{len}\NormalTok{(t1)) }\CommentTok{\# Imprime  la longitud  de la tupla}
\end{Highlighting}
\end{Shaded}

\begin{verbatim}
4
\end{verbatim}

\begin{Shaded}
\begin{Highlighting}[]
\NormalTok{t1[}\BuiltInTok{len}\NormalTok{(t1)}\OperatorTok{{-}}\DecValTok{1}\NormalTok{]  }\CommentTok{\# recupera el último }
\end{Highlighting}
\end{Shaded}

\begin{verbatim}
6
\end{verbatim}
\end{frame}

\begin{frame}[fragile]{Método \texttt{.index(x)}}
\phantomsection\label{muxe9todo-.indexx}
\begin{itemize}
\tightlist
\item
  El método \texttt{.index(x)} regresa el índice que le corresponde al
  elemento \texttt{x} en la tupla, si el elemento está más de una vez
  regresa el índice de la primera ocurrencia, en caso que \texttt{x} no
  esté ocurre un error.
\end{itemize}

\begin{Shaded}
\begin{Highlighting}[]
\NormalTok{t1.index(}\DecValTok{4}\NormalTok{)    }\CommentTok{\# índice correspondiente al primer 4 en t1}
\end{Highlighting}
\end{Shaded}

\begin{verbatim}
1
\end{verbatim}
\end{frame}

\begin{frame}[fragile]{Método \texttt{.count()}}
\phantomsection\label{muxe9todo-.count}
\begin{itemize}
\tightlist
\item
  El método \texttt{.count()} regresa el número de veces que un elemento
  está en la tupla, en caso de no estar regresa cero, a continuación se
  muestra como usarlos sobre la tupla \texttt{t1} creada arriba
\end{itemize}

\begin{Shaded}
\begin{Highlighting}[]
\NormalTok{t1.count(}\DecValTok{4}\NormalTok{) }\CommentTok{\# número de veces que está 4 }
\end{Highlighting}
\end{Shaded}

\begin{verbatim}
1
\end{verbatim}
\end{frame}

\begin{frame}[fragile]{}
\phantomsection\label{section}
\begin{tcolorbox}[enhanced jigsaw, colbacktitle=quarto-callout-important-color!10!white, colframe=quarto-callout-important-color-frame, arc=.35mm, opacityback=0, toptitle=1mm, left=2mm, breakable, title=\textcolor{quarto-callout-important-color}{\faExclamation}\hspace{0.5em}{Importante}, bottomtitle=1mm, titlerule=0mm, rightrule=.15mm, coltitle=black, bottomrule=.15mm, toprule=.15mm, colback=white, leftrule=.75mm, opacitybacktitle=0.6]

Habrá notado la diferencia en el llamado a \texttt{len()} frente a
\texttt{.index()} y \texttt{.count()}:

\begin{itemize}
\item
  \texttt{len()} es una función genérica de Python, así que puede
  aplicarse a diferentes tipos de objetos,
\item
  Mientras que los dos últimos son \emph{métodos}, definidos sobre la
  clase \emph{tuple}.
\end{itemize}

\end{tcolorbox}
\end{frame}

\begin{frame}[fragile]{Diferencia entre método y función}
\phantomsection\label{diferencia-entre-muxe9todo-y-funciuxf3n}
\begin{itemize}
\item
  Cuando se crea una tupla se le asocian los métodos propios de la
  clase, es decir, el método está asociado al objeto.
\item
  Las funciones sólo se pueden llamar por su nombre, ya que se definen
  de forma independiente,
\item
  Los métodos no pueden ser llamados únicamente por su nombre,
  necesitamos invocar la clase por una referencia de esa clase en la que
  está definida.
\item
  Esa es la razón por la que \textbf{no usamos} \texttt{index(t1,4)} o
  \texttt{count(t1,4)} sino \texttt{t1.index(4)} o \texttt{t1.count(4)}
\item
  Porque estos son métodos definidos sobre la clase \texttt{tuple},
  mientras que \texttt{len()} es una función a la cual se le entrega una
  tupla y regresa su longitud,
\item
  \texttt{len()} se puede usar con otros tipos de contenedores, incluso
  con cadenas de caracteres.
\end{itemize}
\end{frame}

\section{Operaciones con tuplas}\label{operaciones-con-tuplas}

\begin{frame}[fragile]{Operador \texttt{in}}
\phantomsection\label{operador-in}
Para saber si un elemento hace parte de una tupla, usamos el operador
\texttt{in} el cual regresa un valor lógico \texttt{True} o
\texttt{False} dependiendo de si el elemento se encuentra o no en la
tupla.

\begin{Shaded}
\begin{Highlighting}[]
\DecValTok{4} \KeywordTok{in}\NormalTok{ t1     }\CommentTok{\# True sí 4 está en t1, False si no  }
\end{Highlighting}
\end{Shaded}

\begin{verbatim}
True
\end{verbatim}

\begin{Shaded}
\begin{Highlighting}[]
\CommentTok{"x"} \KeywordTok{in}\NormalTok{ t1   }\CommentTok{\# True sí  "x" está en t1,  False si no }
\end{Highlighting}
\end{Shaded}

\begin{verbatim}
False
\end{verbatim}

\begin{Shaded}
\begin{Highlighting}[]
\CommentTok{"abc"} \KeywordTok{in}\NormalTok{ t1 }\CommentTok{\# True sí "abc" está en t1,  False si no }
\end{Highlighting}
\end{Shaded}

\begin{verbatim}
False
\end{verbatim}
\end{frame}

\begin{frame}[fragile]{Concatenación de tuplas.}
\phantomsection\label{concatenaciuxf3n-de-tuplas.}
\index{Tuplas!Operaciones con}

Se pueden concatenar y replicar tuplas usando los operadores \texttt{+}
y \texttt{*} de la siguiente forma

\begin{Shaded}
\begin{Highlighting}[]
\NormalTok{Trats}\OperatorTok{=}\NormalTok{(}\StringTok{"T1"}\NormalTok{,}\StringTok{"T2"}\NormalTok{,}\StringTok{"T3"}\NormalTok{) }
\NormalTok{(}\DecValTok{2}\NormalTok{,}\DecValTok{3}\NormalTok{,}\DecValTok{4}\NormalTok{) }\OperatorTok{+}\NormalTok{ Trats }\CommentTok{\# Concatena la tupla (2,3,4) con Trats }
\end{Highlighting}
\end{Shaded}

\begin{verbatim}
(2, 3, 4, 'T1', 'T2', 'T3')
\end{verbatim}

\begin{Shaded}
\begin{Highlighting}[]
\NormalTok{Trats}\OperatorTok{*}\DecValTok{3}         \CommentTok{\# replica la tupla Trats tres veces }
\end{Highlighting}
\end{Shaded}

\begin{verbatim}
('T1', 'T2', 'T3', 'T1', 'T2', 'T3', 'T1', 'T2', 'T3')
\end{verbatim}

\begin{Shaded}
\begin{Highlighting}[]
\NormalTok{(}\DecValTok{2}\NormalTok{,}\DecValTok{3}\NormalTok{,}\DecValTok{4}\NormalTok{)}\OperatorTok{*}\DecValTok{2}       \CommentTok{\# replica la tupla (2,3,4) tres veces }
\end{Highlighting}
\end{Shaded}

\begin{verbatim}
(2, 3, 4, 2, 3, 4)
\end{verbatim}
\end{frame}

\begin{frame}[fragile]{Iterar sobre una tupla.}
\phantomsection\label{iterar-sobre-una-tupla.}
\index{Tuplas!Iterar sobre una}\index{Iterar}

Una tupla es un objeto \emph{iterable}, es decir, se pueden recorrer uno
tras otro sus elementos usando bucles, a continuación un ejemplo usando
\texttt{for}

\begin{Shaded}
\begin{Highlighting}[]
\ControlFlowTok{for}\NormalTok{ i }\KeywordTok{in}\NormalTok{ Trats:}
    \BuiltInTok{print}\NormalTok{(i) }
\end{Highlighting}
\end{Shaded}

\begin{verbatim}
T1
T2
T3
\end{verbatim}
\end{frame}

\begin{frame}[fragile]{Aplicación}
\phantomsection\label{aplicaciuxf3n}
Escribir código que tome una tupla y nos informe cuántos enteros y no
enteros contiene.

\begin{Shaded}
\begin{Highlighting}[]
\NormalTok{t1}\OperatorTok{=}\NormalTok{(}\StringTok{\textquotesingle{}abc\textquotesingle{}}\NormalTok{,}\DecValTok{4}\NormalTok{,}\VariableTok{True}\NormalTok{,}\DecValTok{5}\NormalTok{,}\VariableTok{False}\NormalTok{,}\FloatTok{2.5}\NormalTok{) }
\NormalTok{entero }\OperatorTok{=} \DecValTok{0}
\ControlFlowTok{for}\NormalTok{ i }\KeywordTok{in}\NormalTok{ t1:}
    \ControlFlowTok{if} \BuiltInTok{type}\NormalTok{(i) }\KeywordTok{is} \BuiltInTok{int}\NormalTok{:}
\NormalTok{        entero}\OperatorTok{+=}\DecValTok{1}
\NormalTok{NoEnteros}\OperatorTok{=}\BuiltInTok{len}\NormalTok{(t1)}\OperatorTok{{-}}\NormalTok{entero}
\BuiltInTok{print}\NormalTok{(}\SpecialStringTok{f"Hay }\SpecialCharTok{\{}\NormalTok{entero}\SpecialCharTok{\}}\SpecialStringTok{ enteros en la tupla y }\SpecialCharTok{\{}\NormalTok{NoEnteros}\SpecialCharTok{\}}\SpecialStringTok{ que no son enteros"}\NormalTok{)  }
\end{Highlighting}
\end{Shaded}

\begin{verbatim}
Hay 2 enteros en la tupla y 4 que no son enteros
\end{verbatim}
\end{frame}

\begin{frame}[fragile]{Ejemplo de aplicación.}
\phantomsection\label{ejemplo-de-aplicaciuxf3n.}
En el siguiente bloque de código se generan datos de una población
normal con media cero y desviación estándar 10. Luego se usa la función
\texttt{shapiro} del módulo \texttt{scipy.stats}, para llevar a cabo la
prueba de \emph{Shapiro Wilks} \index{Prueba!de Shapiro Wilks} usando
estos datos. Lo que se quiere resaltar es que la función regresa una
tupla que tiene como primer elemento el estadístico \(W\) y en la
segunda posición el \(p-\)valor de la prueba, veamos
\end{frame}

\begin{frame}[fragile]{Ejemplo de aplicación II}
\phantomsection\label{ejemplo-de-aplicaciuxf3n-ii}
\phantomsection\label{annotated-cell-65}%
\begin{Shaded}
\begin{Highlighting}[numbers=left,,]
\ImportTok{import}\NormalTok{ scipy.stats }\ImportTok{as}\NormalTok{ s }\hspace*{\fill}\NormalTok{\circled{1}}
\ImportTok{import}\NormalTok{ numpy }\ImportTok{as}\NormalTok{ np }\hspace*{\fill}\NormalTok{\circled{2}}
\NormalTok{np.random.seed(seed}\OperatorTok{=}\DecValTok{123}\NormalTok{) }\hspace*{\fill}\NormalTok{\circled{3}}
\NormalTok{n }\OperatorTok{=}\NormalTok{ s.norm(}\DecValTok{0}\NormalTok{,}\DecValTok{10}\NormalTok{) }\hspace*{\fill}\NormalTok{\circled{4}}
\NormalTok{datos}\OperatorTok{=}\NormalTok{n.rvs(}\DecValTok{100}\NormalTok{) }\hspace*{\fill}\NormalTok{\circled{5}}
\NormalTok{res}\OperatorTok{=}\NormalTok{s.shapiro(datos) }\hspace*{\fill}\NormalTok{\circled{6}}
\end{Highlighting}
\end{Shaded}

\begin{description}
\tightlist
\item[\circled{1}]
Carga del modulo \texttt{scipy.stats}.
\item[\circled{2}]
Carga del paquete \texttt{numpy}
\item[\circled{3}]
Establece una semilla para el generador de números aleatorios de
\texttt{NumPy}, para que el ejemplo sea reproducible.
\item[\circled{4}]
Se \texttt{instancia} una variable aleatoria normal de media 0 y
desviación estándar 10.
\item[\circled{5}]
Se simula la toma de una muestra aleatoria de una población con
distribución normal de media 0 y desviación estándar 10.
\item[\circled{6}]
Prueba de Shapiro Wilks.
\end{description}
\end{frame}

\begin{frame}[fragile]{Ejemplo de aplicación III}
\phantomsection\label{ejemplo-de-aplicaciuxf3n-iii}
Podemos recuperar los valores del estadístico de prueba, \(W\), y el
\(p-\)valor, mediante el desempaquetado o mediante el indexado como se
muestra a continuación

\begin{Shaded}
\begin{Highlighting}[]
\NormalTok{W,p}\OperatorTok{=}\NormalTok{res  }\CommentTok{\# desempaquetado, equivalente a W=res[0] ; p=res[1]}
\end{Highlighting}
\end{Shaded}

si estamos escribiendo un informe donde sea necesario reportar el
resultado de la prueba podríamos usar

\begin{Shaded}
\begin{Highlighting}[]
\BuiltInTok{print}\NormalTok{(}\StringTok{"Estadístico de S{-}Wilks, }\SpecialCharTok{\%0.3f}\StringTok{, p\_valor = }\SpecialCharTok{\%0.4f}\StringTok{"} \OperatorTok{\%}\NormalTok{res)}
\end{Highlighting}
\end{Shaded}
\end{frame}

\begin{frame}[fragile]{¿Cómo saber si un objeto de Python es una tupla?}
\phantomsection\label{cuxf3mo-saber-si-un-objeto-de-python-es-una-tupla}
\index{Tuplas!Función \texttt{isinstance()}}
\index{Tuplas!Función \texttt{type()}}

En algunas aplicaciones es necesario verificar si un objeto pertenece o
no a una clase. Por ejemplo, para dirigir el flujo de ejecución del
código de programación de una función: si el objeto entregado es una
tupla, ejecute ciertas líneas de código, de lo contrario ejecute otras
líneas. La función \texttt{isinstance(objeto,tipo)} regresa
\texttt{True} si la clase del \texttt{objeto} entregado corresponde a
\texttt{tipo}, en caso contrario regresa \texttt{False}.

\begin{Shaded}
\begin{Highlighting}[]
\BuiltInTok{isinstance}\NormalTok{(t1, }\BuiltInTok{tuple}\NormalTok{)}
\end{Highlighting}
\end{Shaded}

\begin{verbatim}
True
\end{verbatim}
\end{frame}

\begin{frame}[fragile]{Función \texttt{type()}}
\phantomsection\label{funciuxf3n-type-1}
Otra forma es usando la función \texttt{type()}, como se muestra a
continuación

\begin{Shaded}
\begin{Highlighting}[]
\BuiltInTok{type}\NormalTok{(t1) }\KeywordTok{is} \BuiltInTok{tuple}
\end{Highlighting}
\end{Shaded}

\begin{verbatim}
True
\end{verbatim}

\begin{Shaded}
\begin{Highlighting}[]
\BuiltInTok{type}\NormalTok{(}\DecValTok{4}\NormalTok{) }\KeywordTok{is} \BuiltInTok{tuple}
\end{Highlighting}
\end{Shaded}

\begin{verbatim}
False
\end{verbatim}
\end{frame}

\begin{frame}[fragile]{Otros tipos de datos}
\phantomsection\label{otros-tipos-de-datos}
Esta última línea regresa \texttt{False} porque el número \texttt{4} no
es una tupla. En las siguientes líneas se ilustra como averiguar por los
tipos de datos más usados en estadística

\begin{Shaded}
\begin{Highlighting}[]
\BuiltInTok{type}\NormalTok{(x) }\KeywordTok{is} \BuiltInTok{int}    \CommentTok{\# True si x es entero, False en otro caso}
\BuiltInTok{type}\NormalTok{(x) }\KeywordTok{is} \BuiltInTok{float}  \CommentTok{\# True si x un real, False en otro caso}
\BuiltInTok{type}\NormalTok{(x) }\KeywordTok{is} \BuiltInTok{str}    \CommentTok{\# True si x es string,  False en otro caso}
\BuiltInTok{type}\NormalTok{(x) }\KeywordTok{is} \BuiltInTok{bool}   \CommentTok{\# True si x es lógico,  False en otro caso }
\end{Highlighting}
\end{Shaded}
\end{frame}

\begin{frame}[fragile]{Ejemplo}
\phantomsection\label{ejemplo}
A manera de ejemplo, programar una función que reciba un objeto y en
caso que este sea una tupla regrese el mensaje
\texttt{\textquotesingle{}Es\ una\ tupla\textquotesingle{}}, en caso
contrario regrese
\texttt{\textquotesingle{}No\ es\ una\ tupla\textquotesingle{}}

\begin{Shaded}
\begin{Highlighting}[numbers=left,,]
\CommentTok{\# Se define la función }
\KeywordTok{def}\NormalTok{ verificartupla(x):}
  \ControlFlowTok{if}\NormalTok{(}\BuiltInTok{type}\NormalTok{(x) }\KeywordTok{is} \BuiltInTok{tuple}\NormalTok{):}
    \BuiltInTok{print}\NormalTok{(}\StringTok{"Es una tupla"}\NormalTok{)}
  \ControlFlowTok{else}\NormalTok{:}
    \BuiltInTok{print}\NormalTok{(}\StringTok{"No es una tupla"}\NormalTok{)}
\CommentTok{\# Fin de definición de función}
\end{Highlighting}
\end{Shaded}
\end{frame}

\begin{frame}[fragile]{Uso de la función}
\phantomsection\label{uso-de-la-funciuxf3n}
\begin{Shaded}
\begin{Highlighting}[]
\CommentTok{\# Uso de  la función }
\NormalTok{verificartupla(t1) }\CommentTok{\# Retorna: \textasciigrave{}True\textasciigrave{}, \textasciigrave{}t1\textasciigrave{} es una tupla}
\NormalTok{verificartupla(}\StringTok{"hola"}\NormalTok{) }\CommentTok{\# Retorna: \textasciigrave{}False\textasciigrave{}, \textasciigrave{}"hola"\textasciigrave{} no es una tupla}
\end{Highlighting}
\end{Shaded}
\end{frame}




\end{document}
